% Provenance (never printed): SNT "Les donnees structurees et leur traitement" (descriptors, values, CSV, selection,
% sorting, merging tables, consistency) and NSI premiere "Traitement de donnees en tables" (indexation, search,
% sorting on a column, merging on a common descriptor). See sources/book2/BRIEFS.md.
\chapter{Processing Data Tables}\label{ch:g10:data-tables}

The town of Brenmoor has two weather stations, one on the hill and one by the river, and each records the
temperature, the rain and the wind once a day. Ilias wants the year's mean temperature at the river for a school
report. His first try gives $11.43$~°C. Then he notices a January day at $-99$~°C, a summer day recorded twice, and
blank cells. Which number is true?

This chapter treats a file of records as a table and does to it what every data program does: select, sort, group,
\omterm{def:g10:data-tables:join}{join} with another table, and find the records that cannot be right. Each operation is a few lines of Python; the
care goes into knowing what each line does to the data.

\begin{recall}
A dictionary maps keys to values; a record can be stored as a dictionary, and a table as a list of records. A table
saved as a text file with separators can be loaded as such a list (Book 1, grade 8, which also sorted lists and
counted comparisons). The mathematics book defines the mean, the median and the spread of a list of values
(grade 10).
\end{recall}

\section{A table of records}\label{sec:g10:data-tables:records}

\begin{definition}[Header, descriptor, domain]\label{def:g10:data-tables:header}
A data table is a list of records that all have the same fields. The name of a field is a
\emph{descriptor}\index{descriptor}; the first line of a table file, which lists the descriptors, is its
\emph{header}\index{header (of a table)}. The \emph{domain}\index{domain of a column} of a column is the set of values
it may take: dates of 2025, the two station names, temperatures in degrees.
\end{definition}

\begin{omfigure}
\begin{tikzpicture}[x=24mm, y=-6mm, font=\footnotesize\ttfamily]
  \fill[omDef!12] (0,0) rectangle (5,1);
  \fill[omProp!14] (0,2) rectangle (5,3);
  \fill[omMeth!14] (2,0) rectangle (3,4);
  \foreach \c in {0,...,5} \draw[black!50] (\c,0) -- (\c,4);
  \foreach \r in {0,...,4} \draw[black!50] (0,\r) -- (5,\r);
  \foreach \h [count=\c from 0] in {date,station,temp\_c,rain\_mm,wind\_kmh}
    \node at (\c+0.5,0.5) {\h};
  \foreach \r/\a/\b/\t/\u/\v in {1/2025-01-01/Hilltop/3.9/0.0/15.3, 2/2025-01-01/Riverside/4.3/0.0/10.4,
      3/2025-01-02/Hilltop/3.9/0.0/12.4} {
    \node at (0.5,\r+0.5) {\a}; \node at (1.5,\r+0.5) {\b}; \node at (2.5,\r+0.5) {\t};
    \node at (3.5,\r+0.5) {\u}; \node at (4.5,\r+0.5) {\v};
  }
  \draw[very thick, omThm] (4,1) rectangle (5,2);
  \begin{scope}[font=\footnotesize\rmfamily]
    \node[anchor=east, text=omDef] at (-0.05,0.5) {header};
    \node[anchor=east, text=omProp] at (-0.05,2.5) {a record (row)};
    \node[anchor=north, text=omMeth] at (2.5,4.1) {a column};
    \node[anchor=west, text=omThm] at (5.05,1.5) {a cell};
  \end{scope}
\end{tikzpicture}

{\small The first lines of the weather table: a \omterm{def:g10:data-tables:header}{header} of five \omterm{def:g10:data-tables:header}{descriptors}, then one record per station and per
day.}
\end{omfigure}

% ledger: g10:csv-4180
The town's tables are \emph{CSV} files (comma-separated values, described in 2005): one record per line, the fields
separated by commas, the \omterm{def:g10:data-tables:header}{header} first. A field that itself holds a comma is put between double quotes. Python's
\code{csv} module knows these rules; \code{csv.DictReader} reads each line as a dictionary whose keys are the
\omterm{def:g10:data-tables:header}{descriptors}. Every value it gives is a string, so the program converts them.

\omcode{code/grade-10/07-data-tables/listings/load_weather.py}{1}{30}{Loading a CSV table with converted
values.}\label{lst:g10:data-tables:load-weather}

\omoutput{code/grade-10/07-data-tables/outputs/load_weather.txt}{$731$ records, but a year has $365$ days: Riverside
has one record too many.}\label{out:g10:data-tables:load-weather}

The other listings of the chapter import \code{read_table} from this file. An empty cell becomes \code{None}, a
value that Python reserves for ``nothing''; keeping it apart from $0$ matters, since a day without a measured
temperature is not a day at $0$~°C.

\section{Selecting, projecting and sorting}\label{sec:g10:data-tables:select}

\begin{definition}[Selection, projection]\label{def:g10:data-tables:selection}
A \emph{selection}\index{selection} keeps the records of a table that satisfy a condition. A
\emph{projection}\index{projection} keeps some columns of every record.
\end{definition}

\begin{omfigure}
\begin{tikzpicture}[x=7mm, y=-4mm]
  \foreach \r in {1,4,5} \fill[omProp!30] (0,\r) rectangle (5,\r+1);
  \foreach \r in {0,...,6} \foreach \c in {0,...,4} \draw[black!40] (\c,\r) rectangle (\c+1,\r+1);
  \node[font=\footnotesize, below] at (2.5,7.2) {selection: some rows};
  \begin{scope}[xshift=55mm]
    \foreach \c in {0,3} \fill[omMeth!30] (\c,0) rectangle (\c+1,7);
    \foreach \r in {0,...,6} \foreach \c in {0,...,4} \draw[black!40] (\c,\r) rectangle (\c+1,\r+1);
    \node[font=\footnotesize, below] at (2.5,7.2) {projection: some columns};
  \end{scope}
\end{tikzpicture}

{\small \omterm{def:g10:data-tables:selection}{Selection} keeps whole rows that pass a test; \omterm{def:g10:data-tables:selection}{projection} keeps whole columns. Both give a new table.}
\end{omfigure}

\omcode{code/grade-10/07-data-tables/listings/select_project.py}{1}{9}{The rainiest days on the hill: a \omterm{def:g10:data-tables:selection}{selection},
then a \omterm{def:g10:data-tables:selection}{projection}.}\label{lst:g10:data-tables:select-project}

\omoutput{code/grade-10/07-data-tables/outputs/select_project.txt}{Four days of $15$~mm of rain or more at
Hilltop.}\label{out:g10:data-tables:select-project}

Sorting a table needs a rule saying which record comes first.

\begin{definition}[Sort key, stable sort]\label{def:g10:data-tables:sort-key}
A \emph{sort key}\index{sort key} is a function that gives each record the value it is sorted by. A sort is
\emph{stable}\index{stable sort} when records with equal keys keep the order they had before sorting.
\end{definition}

Python's \code{sorted(table, key=f)} sorts by \code{f(record)}, and it is \omterm{def:g10:data-tables:sort-key}{stable}. A key may be a tuple: tuples
compare their first items, then, on a tie, their second items, and so on. A minus sign in front of a number sorts it
from largest to smallest.

\omcode{code/grade-10/07-data-tables/listings/sort_key.py}{1}{11}{Sorting on one key, then on two
keys.}\label{lst:g10:data-tables:sort-key}

\omoutput{code/grade-10/07-data-tables/outputs/sort_key.txt}{The four windiest days; then the first four records of
10 to 19 June, sorted by station and, within a station, by rain from most to least.}\label{out:g10:data-tables:sort-key}

\begin{remark}
Since \code{sorted} is \omterm{def:g10:data-tables:sort-key}{stable}, the same order can be built in two passes: first by rain, then by station. The second
pass keeps the rain order among the records of each station. \Cref{exo:g10:data-tables:13} proves it.
\end{remark}

\section{Grouping and aggregating}\label{sec:g10:data-tables:groups}

\begin{definition}[Grouping, aggregate]\label{def:g10:data-tables:aggregate}
\emph{Grouping}\index{grouping} splits the records of a table into groups that share the value of one or more
columns. An \emph{aggregate}\index{aggregate} is one value computed from a whole group: a count, a sum, a mean, a
median, a largest value.
\end{definition}

The usual way in Python is a dictionary from the group's key to the list of its values, filled in one pass.

\omcode{code/grade-10/07-data-tables/listings/group_mean.py}{1}{14}{Monthly mean temperature per
station.}\label{lst:g10:data-tables:group-mean}

\omoutput{code/grade-10/07-data-tables/outputs/group_mean.txt}{Five of the twelve months. Riverside is warmer than
Hilltop, except in January and December.}\label{out:g10:data-tables:group-mean}

The river station is lower and warmer all summer, yet colder than the hill in January and December by about two
degrees. Before believing it, look at what went into those two means: each holds one of the $-99$ values. One
reading of $-99$ among $31$ lowers a monthly mean by about $(99 + 3)/31 \approx 3.3$ degrees.

\section{Joining tables}\label{sec:g10:data-tables:join}

The library keeps its loans in one table and its books in another; a loan names its book only by
\code{book_id}. To print the titles of the books still out, both tables are needed.

\begin{definition}[Join]\label{def:g10:data-tables:join}
The \emph{join}\index{join} of two tables on a shared column is the table of the pairs of records, one from each
table, that have the same value in that column, each pair merged into one record.
\end{definition}

\begin{omfigure}
\begin{tikzpicture}[x=17mm, y=-6mm, font=\footnotesize]
  \node[font=\footnotesize\bfseries] at (1,-0.6) {loans};
  \fill[omDef!12] (0,0) rectangle (2,1);
  \node at (0.5,0.5) {\texttt{loan\_id}}; \node at (1.5,0.5) {\texttt{book\_id}};
  \foreach \r/\l/\b in {1/3/28, 2/11/28, 3/27/12} {
    \node at (0.5,\r+0.5) {\l}; \node (L\r) at (1.5,\r+0.5) {\b};
  }
  \foreach \c in {0,1,2} \draw[black!40] (\c,0) -- (\c,4);
  \foreach \r in {0,...,4} \draw[black!40] (0,\r) -- (2,\r);
  \begin{scope}[xshift=60mm]
    \node[font=\footnotesize\bfseries] at (1.25,-0.6) {books};
    \fill[omDef!12] (0,0) rectangle (2.5,1);
    \node at (0.5,0.5) {\texttt{book\_id}}; \node at (1.75,0.5) {\texttt{title}};
    \foreach \r/\b/\t in {1/12/{Under the Forest}, 2/28/{Under the Bridge}, 3/40/{The Lost Code}} {
      \node (B\r) at (0.5,\r+0.5) {\b}; \node at (1.75,\r+0.5) {\t};
    }
    \foreach \c in {0,1,2.5} \draw[black!40] (\c,0) -- (\c,4);
    \foreach \r in {0,...,4} \draw[black!40] (0,\r) -- (2.5,\r);
  \end{scope}
  \draw[omProp, thick] (2,1.5) -- ([xshift=60mm]0,2.5);
  \draw[omProp, thick] (2,2.5) -- ([xshift=60mm]0,2.5);
  \draw[omProp, thick] (2,3.5) -- ([xshift=60mm]0,1.5);
\end{tikzpicture}

{\small Joining on \texttt{book\_id}, on three loans and three books of the town's tables: each loan meets the book
with its number. Book $28$ appears in two joined records; book $40$, which none of these three loans names, appears in
none.}
\end{omfigure}

The simplest program compares every record of one table with every record of the other: a \emph{nested-loop \omterm{def:g10:data-tables:join}{join}}.

\omcode{code/grade-10/07-data-tables/listings/join_loans.py}{1}{16}{A nested-loop \omterm{def:g10:data-tables:join}{join} of the loans still out with
the books.}\label{lst:g10:data-tables:join-loans}

\omoutput{code/grade-10/07-data-tables/outputs/join_loans.txt}{$22 \times 80 = 1760$ comparisons for $22$
matches.}\label{out:g10:data-tables:join-loans}

A dictionary does better: build it once from the books, by number, and each loan then finds its book in one lookup.
The listing \omterm{def:g10:data-tables:join}{joins} four tables this way, from the loan to the book, and from the loan to the member and to the person.

\omcode{code/grade-10/07-data-tables/listings/join_index.py}{1}{17}{Joining through dictionary indexes: overdue loans
with the title and the borrower.}\label{lst:g10:data-tables:join-index}

\omoutput{code/grade-10/07-data-tables/outputs/join_index.txt}{The ten loans still out that were due before
July.}\label{out:g10:data-tables:join-index}

\begin{proposition}[Cost of a join]\label{prop:g10:data-tables:join-cost}
Joining a table of $n$ records with a table of $m$ records takes $n \times m$ comparisons with nested loops; with a
dictionary built on the second table, it takes $m$ insertions and $n$ lookups.
\end{proposition}

For the $400$ loans and $80$ books, that is $32{,}000$ comparisons against $480$ steps. The cost of an algorithm is
the subject of grade 11; here it is enough to count.

\section{Missing and inconsistent data}\label{sec:g10:data-tables:clean}

\begin{definition}[Missing value, consistency test]\label{def:g10:data-tables:missing}
A \emph{missing value}\index{missing value} is a cell with no data: empty in the file, \code{None} in the program.
A \emph{consistency test}\index{consistency test} is a rule that every correct record satisfies, such as a value
inside its column's \omterm{def:g10:data-tables:header}{domain}, or a due date $21$ days after the loan; a record that fails it is wrong, or the rule is.
\end{definition}

Error codes are the worst kind of bad value: a station that cannot measure writes $-99.0$, a number inside the
column's type but far outside its \omterm{def:g10:data-tables:header}{domain}, and every sum and mean takes it at face value.

\omcode{code/grade-10/07-data-tables/listings/clean_weather.py}{1}{28}{Finding the defects, and the means with and
without them.}\label{lst:g10:data-tables:clean-weather}

\omoutput{code/grade-10/07-data-tables/outputs/clean_weather.txt}{Six \omterm{def:g10:data-tables:missing}{missing values}, two error codes, one duplicate
record; the Riverside mean moves by more than half a degree.}\label{out:g10:data-tables:clean-weather}

\begin{omfigure}
\begin{tikzpicture}
  \begin{axis}[omrecord, width=0.85\linewidth, height=5.5cm, xlabel={month of 2025}, ylabel={mean temperature (°C)},
      xmin=0.5, xmax=12.5, xtick={1,...,12}, ymin=-1, ymax=23, legend pos=north west, legend cell align=left]
    \addplot[omThm, thick, dashed, mark=o] table[x=month, y=riverside_raw]
      {code/grade-10/07-data-tables/figures/monthly_means.dat};
    \addlegendentry{Riverside, raw}
    \addplot[omDef, thick, mark=*] table[x=month, y=riverside]
      {code/grade-10/07-data-tables/figures/monthly_means.dat};
    \addlegendentry{Riverside, cleaned}
    \addplot[omMeth, thick, mark=square*] table[x=month, y=hilltop]
      {code/grade-10/07-data-tables/figures/monthly_means.dat};
    \addlegendentry{Hilltop, cleaned}
  \end{axis}
\end{tikzpicture}

{\small Monthly means of the two stations. The two error codes pull January and December of the raw river record
below the hill; once they are removed, the river station is the warmer one in every month.}
\end{omfigure}

\begin{method}[Cleaning a table]\label{met:g10:data-tables:clean}
\begin{enumerate}
  \item Count the records and compare with what is expected (one per station and per day).
  \item For each column, check every value against the \omterm{def:g10:data-tables:header}{domain}; list the \omterm{def:g10:data-tables:missing}{missing values} and the impossible ones.
  \item Look for records that repeat a key that should be unique, such as (date, station).
  \item Decide, for each kind of defect, a rule: drop the record, keep it without that value, or correct it from
    another source, and write the rule down next to the result.
  \item Never change the original file: write the cleaned table to a new one, or clean while reading.
\end{enumerate}
\end{method}

\begin{inpractice}[Data about people]
The library tables are about people: who borrowed what, and when. A table that lets someone be identified,
directly by a name or indirectly by a joined number, holds personal data, which data protection law lets an
organisation keep only for a stated purpose, and for no longer than that purpose needs. A \omterm{def:g10:data-tables:join}{join} is exactly how
scattered harmless columns become a profile of a person.
\end{inpractice}

Not every table is flat. A member has several loans, and a loan list inside a record does not fit in a CSV line.

\begin{definition}[JSON]\label{def:g10:data-tables:json}
% ledger: g10:json-8259
\emph{JSON}\index{JSON} (JavaScript Object Notation, specified in its current form in 2017) is a text format for
nested data: objects between braces, which map names to values, lists between square brackets, and numbers,
strings, \code{true}, \code{false} and \code{null}. A value may itself be an object or a list. JSON exchanged
between systems is written in \omterm{def:g10:text-encodings:utf8}{UTF-8}.
\end{definition}

\omcode{code/grade-10/07-data-tables/listings/to_json.py}{1}{13}{A member and their first loans as nested
\omterm{def:g10:data-tables:json}{JSON}.}\label{lst:g10:data-tables:to-json}

\omoutput{code/grade-10/07-data-tables/outputs/to_json.txt}{\code{None} is written \code{null}, and reading the text
back gives the same dictionary.}\label{out:g10:data-tables:to-json}

Python's \code{json.dumps} turns dictionaries and lists into \omterm{def:g10:data-tables:json}{JSON} text, and \code{json.loads} reads it back; web
services exchange their data this way, as the next chapter shows.

\section{Exercises}\label{sec:g10:data-tables:exercises}

\begin{exercise}[$\star$]\label{exo:g10:data-tables:1}
Change the condition of the \omterm{def:g10:data-tables:selection}{selection} listing to \code{row["rain_mm"] >= 15} without the station test. Will the
result have more or fewer records than before? Can it have fewer than $4$?
\end{exercise}

\begin{exercise}[$\star$]\label{exo:g10:data-tables:2}
What does \code{sorted(["pear", "fig", "apple", "kiwi"], key=len)} return? Why does \code{"kiwi"} stay after
\code{"pear"}?
\end{exercise}

\begin{exercise}[$\star$]\label{exo:g10:data-tables:3}
Write \code{project(table, cols)}, which returns the table restricted to the columns of the list \code{cols}, in that
order.
\end{exercise}

\begin{exercise}[$\star$]\label{exo:g10:data-tables:4}
Table $A$ has $5$ records whose key values are $1, 1, 2, 3, 3$; table $B$ has $4$ records with keys $1, 2, 2, 4$. How
many records does the \omterm{def:g10:data-tables:join}{join} of $A$ and $B$ on the key have? Which keys never appear in it?
\end{exercise}

\begin{exercise}[$\star$]\label{exo:g10:data-tables:5}
A program reads the books with \code{csv.DictReader} and sorts them with \code{key=lambda r: r["year"]}, without
converting anything. Why does a book of year \code{10} come before one of year \code{9}? Correct it. Why does the
mistake not show on the town's books, whose years all have four digits?
\end{exercise}

\begin{exercise}[$\star\star$]\label{exo:g10:data-tables:6}
Write \code{group_count(table, col)}, which returns a dictionary from each value of the column to the number of
records with that value. Use it to count the books of each genre.
\end{exercise}

\begin{exercise}[$\star\star$]\label{exo:g10:data-tables:7}
A duplicate detector stores each line of a file in a set and reports a line already seen. The file holds
\code{2025-06-14,Riverside,21.5} and, further on, the same record with a space after the first comma. Why is the
duplicate missed? Correct the detector.
\end{exercise}

\begin{exercise}[$\star\star$]\label{exo:g10:data-tables:8}
Records $(\text{name}, \text{class})$: (Ada, B), (Ben, A), (Cy, B), (Dee, A), (Eli, B). Give the result of a \omterm{def:g10:data-tables:sort-key}{stable
sort} on the class; then of a \omterm{def:g10:data-tables:sort-key}{stable sort} on the name followed by a \omterm{def:g10:data-tables:sort-key}{stable sort} on the class.
\end{exercise}

\begin{exercise}[$\star\star$]\label{exo:g10:data-tables:9}
Write \code{median_by(table, group, col)}, which returns the median of the column \code{col} for each value of
\code{group}, skipping \omterm{def:g10:data-tables:missing}{missing values}. Compare the median and the mean temperature of the raw Riverside record: which
is less disturbed by the error codes, and why?
\end{exercise}

\begin{exercise}[$\star\star$]\label{exo:g10:data-tables:10}
How many comparisons does a nested-loop \omterm{def:g10:data-tables:join}{join} of the $400$ loans with the $80$ books make, and how many steps does a
\omterm{def:g10:data-tables:join}{join} through a dictionary make? And for $100{,}000$ loans and $20{,}000$ books?
\end{exercise}

\begin{exercise}[$\star\star$]\label{exo:g10:data-tables:11}
A program computes the mean temperature by skipping \code{None} values only. On Riverside it prints $11.43$. Explain
what is wrong, and correct the function so that it also skips a given list of error codes.
\end{exercise}

\begin{exercise}[$\star\star\star$]\label{exo:g10:data-tables:12}
Write \code{left_join(left, right, key)}: like a \omterm{def:g10:data-tables:join}{join}, but a record of \code{left} without any match is kept once,
with \code{None} for the columns of \code{right}. Use it to find the members who borrowed nothing in 2025, and the
books that were never lent. How many records does the left \omterm{def:g10:data-tables:join}{join} of the members with the loans have?
\end{exercise}

\begin{exercise}[$\star\star\star$]\label{exo:g10:data-tables:13}
Prove that sorting a table stably by a key $k_2$, then stably by a key $k_1$, gives the table sorted by the pair
$(k_1, k_2)$. Take two records $r$ and $s$ with $(k_1(r), k_2(r)) < (k_1(s), k_2(s))$ and show that $r$ ends before
$s$, in each of the two possible cases.
\end{exercise}

\begin{exercise}[$\star\star\star$]\label{exo:g10:data-tables:14}
Write a consistency checker for the loans: every due date is $21$ days after its out date, and no book comes back
before it went out. Use \code{datetime.date.fromisoformat} and \code{timedelta}. Run it on the town's loans; then
change one line of a copy and check that it is caught.
\end{exercise}

\begin{exercise}[$\star\star\star$]\label{exo:g10:data-tables:15}
Export all members with their list of loans as \omterm{def:g10:data-tables:json}{JSON}, read the text back, and test that the result equals what was
exported. Which Python values would not survive such a round trip?
\end{exercise}

\section{Problem: The Corrected Weather Record}\label{sec:g10:data-tables:problem}

\begin{problem}[{Weekend problem --- what the river station really measured in 2025}]\label{pb:g10:data-tables:1}
The program \code{exercises/corrected_weather.py} cleans the weather table and computes the year's mean temperature
and warmest month of each station. Use it and the listings of the chapter.

\needspace{4\baselineskip}\medskip\noindent\textbf{Part I --- Loading and counting.}
\begin{enumerate}
  \item How many records should the table hold for two stations over the year 2025? How many does it hold?
  \item How many records does each station have?
  \item Which \omterm{def:g10:data-tables:header}{descriptors} does the \omterm{def:g10:data-tables:header}{header} give, and what is the \omterm{def:g10:data-tables:header}{domain} of each column?
  \item Which value does a blank temperature cell become after \code{read_table}?
\end{enumerate}

\needspace{4\baselineskip}\medskip\noindent\textbf{Part II --- Finding the defects.}
\begin{enumerate}[resume]
  \item List the defects that \code{clean_weather.py} finds, by kind and by station.
  \item Which \omterm{def:g10:data-tables:missing}{consistency test} reveals each kind of defect?
  \item Are the two records of 14 June at Riverside identical? What does that suggest about their origin?
  \item How many temperature values does each station have once \omterm{def:g10:data-tables:missing}{missing values} are left out?
\end{enumerate}

\needspace{4\baselineskip}\medskip\noindent\textbf{Part III --- Rules.}
\begin{enumerate}[resume]
  \item For a missing temperature, would you drop the whole record or keep its rain and wind? Why?
  \item For a record with $-99.0$, why must the value not be replaced by $0$?
  \item For the duplicate, which record do you keep?
  \item Could a \omterm{def:g10:data-tables:missing}{missing value} be filled from the other station on the same day? What would that assume?
  \item How many records remain once the rules of the program are applied?
\end{enumerate}

\needspace{4\baselineskip}\medskip\noindent\textbf{Part IV --- The corrected record.}
\begin{enumerate}[resume]
  \item Run the program. Give the raw and the corrected mean temperature of each station.
  \item Why does the Hilltop mean not change?
  \item What would the Riverside mean be with only the duplicate removed, the codes still in? Compare with the
    corrected one.
  \item Check the size of the codes' effect by hand: two values of $-99$ added to $360$ good values whose mean is
    $12.016$, making $362$.
  \item Which month is the warmest at each station, and by how much is the river warmer than the hill then?
  \item State the result: the corrected 2025 mean temperature at Riverside, the number of days it rests on, and the
    error that the error codes would have caused.
\end{enumerate}
\end{problem}
